% Chapter 6 -- Claim: High-Interest Currencies Depreciate Enough to Offset Their Yield
% Requires: amsmath, amssymb, booktabs, array, tabularx, longtable, graphicx, tcolorbox (breakable, skins), hyperref or url.
% Defines the callout box only if the preamble has not.
\providecolor{fenavy}{HTML}{1B2A6B}
\providecolor{feshade}{HTML}{F3F5FA}
\makeatletter
\@ifundefined{callout}{%
  \newtcolorbox{callout}{breakable, enhanced,
    colback=white, colframe=gray!70, boxrule=0.5pt, arc=1mm,
    borderline west={2.5pt}{0pt}{gray!70},
    left=8pt, right=8pt, top=4pt, bottom=4pt}%
}{}
\makeatother

\chapter{Claim: High-Interest Currencies Depreciate Enough to Offset Their Yield}
\label{ch:06}

When the Bretton Woods system of fixed exchange rates ended in 1973, currencies began to float, and economists gained a new laboratory. Suppose investors care only about expected returns. Then a currency that pays a higher interest rate must be expected to fall by just enough to cancel its extra yield. Otherwise everyone would borrow in the low-rate currency and lend in the high-rate one until the opportunity disappeared. This idea is called \textbf{uncovered interest parity}\index{uncovered interest parity (UIP)|textbf}, and much of international macroeconomics rests on it. Fama (1984) tested it with the floating exchange rates of the 1970s. Instead of finding that high-interest currencies fell, as the theory predicts, he found that they tended to rise.

This result became known as the \textbf{forward premium puzzle}\index{forward premium puzzle}, one of the most persistent findings in finance. It has a practical counterpart in the \textbf{carry trade}\index{carry trade}: borrowing in a low-interest currency to invest in a high-interest one. In 2007 a Japanese household could borrow yen at less than 1\% and deposit Australian dollars at more than 6\%. Japanese retail investors, nicknamed ``Mrs.\ Watanabe,'' did so on a large scale, alongside hedge funds. Later research asked why the trade pays. Brunnermeier, Nagel and Pedersen (2008) showed that it suffers sudden crashes when the borrowed currency, such as the yen, jumps in value. Lustig, Roussanov and Verdelhan (2011) grouped currencies by their interest rates and found a common risk that explains much of the return.

The answer matters to central banks, companies and investors. If interest parity fails, a central bank that raises rates may see its currency rise rather than fall. A company that leaves its foreign revenue unhedged is taking a predictable bet. And an investor in the carry trade is being paid to bear crash risk, not collecting a free lunch. The claim on trial\index{claims-on-trial method} is the textbook one: \textbf{high-interest currencies fall enough to cancel their extra yield, so the carry trade earns nothing on average.}

The lab in Section~\ref{sec:06-lab} repeats Fama's (1984) test for nine currencies against the dollar, using free exchange rates and interest rates from FRED. It then builds the carry portfolios of Lustig, Roussanov and Verdelhan (2011). The goal is to check whether uncovered interest parity holds today, and what the carry trade's returns and crashes look like.

By the end of the chapter you will be able to:

\begin{itemize}
  \item compute a forward exchange rate from covered interest parity\index{covered interest parity (CIP)}, and the profit or loss on a carry trade;
  \item run and interpret the Fama regression, and explain the forward premium puzzle\index{forward premium puzzle};
  \item compare the explanations for the puzzle and the evidence for each;
  \item judge what a currency carry fund is paid for and how it loses money.
\end{itemize}

\section{Foundations: exchange rates, forwards and interest parity}
\label{sec:06-foundations}

\subsection{Exchange rates}

The \textbf{spot rate}\index{spot rate|textbf} $S$ is the price of one currency in terms of another. In this chapter $S$ is the number of U.S. dollars per unit of foreign currency, so a rise in $S$ means the foreign currency has \textbf{appreciated}, or become more valuable. Changes are measured in logs: $\Delta s_{t+1} = \ln(S_{t+1}/S_t)$.

\subsection{Forward rates and covered interest parity (CIP)}

A forward contract fixes today the exchange rate $F$ for a trade at a future date. An investor with dollars has two safe ways to earn interest. She can earn the dollar rate $i$ directly. Or she can convert her dollars into foreign currency, earn the foreign rate $i^*$, and lock in the rate for converting back with a forward\index{forward rate} contract. Both routes are free of currency risk, so arbitrage\index{arbitrage} makes them pay the same:
\begin{equation}
  \frac{F}{S} = \frac{1 + i}{1 + i^*}, \qquad \text{or in logs,} \qquad f - s \approx i - i^*
  \label{eq:06-cip}
\end{equation}
For a high-interest currency, the forward rate is therefore below the spot rate. The gap, called the \textbf{forward discount}\index{forward discount|textbf}, equals the difference in interest rates. Covered interest parity held closely until 2008, and has shown small but persistent deviations since (Du, Tepper and Verdelhan 2018).

\subsection{Uncovered interest parity (UIP)}

\textbf{Uncovered interest parity (UIP)}\index{uncovered interest parity (UIP)|textbf} goes one step further. It says that the spot rate \emph{expected} in the future equals today's forward rate, so a high-interest currency is expected to fall by its interest advantage:
\begin{equation}
  E_t[\Delta s_{t+1}] = i_t - i^*_t
  \label{eq:06-uip}
\end{equation}
Unlike covered interest parity, UIP is not enforced by arbitrage. Holding a foreign currency without a forward contract is risky, so UIP holds only if investors demand no extra return for that risk.

\subsection{The carry trade}

The carry trade borrows in a low-interest currency, converts the money into a high-interest currency, invests it there, and converts it back later without any hedge. Its excess return\index{return!excess} is the interest gap plus the change in the exchange rate:
\begin{equation}
  rx_{t+1} = (i^*_t - i_t) + \Delta s_{t+1}
  \label{eq:06-carry}
\end{equation}
Under UIP its expected value is zero, because the currency is expected to fall by exactly the interest gap. Put another way, the carry trade buys high-interest currencies at a discount in the forward market and hopes they will not fall by as much as the discount.

\textbf{Worked example.} The U.S. one-year rate is 1\%, the Australian rate is 5\%, and the spot rate is \$0.70 per Australian dollar (Figure~\ref{fig:06-carry}).

\begin{itemize}
  \item Covered interest parity sets the one-year forward rate at $0.70 \times 1.01/1.05 \approx \$0.673$.
  \item Under UIP, the spot rate expected in a year is also \$0.673, a fall of 3.8\%.
  \item The carry trade borrows \$1 million at 1\%, buys A\$1.429 million, and invests it at 5\% to receive A\$1.5 million in a year. Its result depends on the spot rate then:
  \begin{itemize}
    \item if the rate is unchanged at \$0.70, the proceeds are \$1.05 million, the loan costs \$1.01 million, and the profit is \$40,000;
    \item if the rate falls to \$0.673, as UIP predicts, the proceeds are \$1.01 million and the profit is zero;
    \item if the rate falls 10\% to \$0.63, the proceeds are \$945,000 and the loss is \$65,000.
  \end{itemize}
\end{itemize}

\begin{figure}[htbp]
\centering
\scalebox{0.8}{%
\begin{tikzpicture}[
  % Profit ($ thousands) = 1500 x (spot - 0.6733): +40k if unchanged at 0.70,
  % zero at the UIP forecast 0.673, -65k after a 10% fall to 0.63
  declare function={carry(\s) = 1500*(\s - 0.67333);}
]
\begin{axis}[
  width=0.61\linewidth, height=0.34\linewidth,
  axis lines=left, axis line style={-stealth},
  xmin=0.6, xmax=0.75, ymin=-110, ymax=120,
  xtick={0.6,0.62,...,0.74}, ytick={-100,0,100},
  xticklabel style={/pgf/number format/fixed, /pgf/number format/precision=2},
  xlabel={Spot rate in one year (\$ per Australian dollar)},
  ylabel={Carry profit (\$ thousands)},
  clip=false,
]
  \addplot[gray, thin, domain=0.6:0.75, samples=2] {0};
  \addplot[fenavy, very thick, domain=0.6:0.75, samples=2] {carry(x)};
  \fill[fenavy] (axis cs:0.70,40)  circle (2.6pt);
  \fill[fenavy] (axis cs:0.67333,0) circle (2.6pt);
  \fill[fenavy] (axis cs:0.63,-65) circle (2.6pt);
  \node[anchor=south east] at (axis cs:0.697,42) {unchanged: $+\$40$k};
  \node[anchor=north west] at (axis cs:0.676,-3) {UIP forecast 0.673: zero};
  \node[anchor=west] at (axis cs:0.633,-70) {$-10\%$: $-\$65$k};
\end{axis}
\end{tikzpicture}%
}
\caption{Profit on the worked example's carry trade against the spot rate a year later. The trade earns \$40,000 if the Australian dollar is unchanged, breaks even at the UIP forecast of \$0.673, and loses \$65,000 if the currency falls 10\%.}
\label{fig:06-carry}
\end{figure}

\section{How it works}
\label{sec:06-how}

UIP can be tested in two ways: one currency at a time, or across many currencies at once.

\textbf{The Fama regression.} For each currency, regress the change in the exchange rate on the interest gap at the start of the period:
\begin{equation}
  \Delta s_{t+1} = a + b\,(i_t - i^*_t) + e_{t+1}
  \label{eq:06-fama}
\end{equation}
UIP implies $b = 1$: a foreign currency with a one-point interest advantage should fall by one point. If $b = 0$, exchange rates move at random, and the carry trade keeps the whole interest gap. If $b < 0$, high-interest currencies tend to \emph{rise}, and the carry trade earns even more than the interest gap (Figure~\ref{fig:06-fama}). This last case is the \textbf{forward premium puzzle}.

\begin{figure}[htbp]
\centering
\scalebox{0.8}{%
\begin{tikzpicture}
\begin{axis}[
  width=0.46\linewidth, height=0.34\linewidth,
  axis lines=left, axis line style={-stealth},
  xmin=-0.5, xmax=4.5, ymin=-4.5, ymax=4.5,
  xtick={0,1,...,4}, ytick={-4,-2,0,2,4},
  xlabel={Foreign interest advantage, $i^* - i$ (\%)},
  ylabel={Expected change in foreign currency (\%)},
  clip=false,
]
  \addplot[gray, thin, domain=-0.5:4.5, samples=2] {0};
  % Typical estimate b = -0.9: the currency rises by 0.9 x its advantage
  \addplot[thick, densely dotted, domain=0:4.3, samples=2] {0.9*x};
  % Random walk b = 0
  \addplot[fegray, thick, dashed, domain=0:4.3, samples=2] {0};
  % UIP b = 1: falls by its interest advantage
  \addplot[fenavy, very thick, domain=0:4.3, samples=2] {-x};
  \node[anchor=west, align=left] at (axis cs:4.6,3.87)
    {typical estimate\\($b \approx -0.9$):\\rises instead};
  \node[fegray, anchor=west, align=left] at (axis cs:4.6,0)
    {random walk\\($b = 0$)};
  \node[fenavy, anchor=west, align=left] at (axis cs:4.6,-4.3)
    {UIP ($b = 1$):\\falls by its\\interest advantage};
\end{axis}
\end{tikzpicture}%
}
\caption{Three readings of the Fama regression. Under UIP ($b = 1$), a currency with a higher interest rate falls by its interest advantage. Under a random walk ($b = 0$), it is expected not to move. The typical estimate ($b \approx -0.9$, Froot and Thaler 1990) implies that it rises instead.}
\label{fig:06-fama}
\end{figure}

\textbf{Carry portfolios.} Each month, rank currencies by their interest rates, buy the highest-yielding ones and sell the lowest-yielding ones. If UIP holds, this portfolio earns nothing on average. A portfolio averages out the noise of individual currencies, and it is what investors actually trade.

\textbf{Why UIP might fail.} There are three broad explanations.

\begin{itemize}
  \item \textbf{A risk premium.} High-interest currencies tend to fall in global downturns, when investors flee to the dollar, the yen and the Swiss franc. Investors therefore demand extra return for holding them. Lustig, Roussanov and Verdelhan (2011) argue that carry returns are payment for exposure to this global risk.
  \item \textbf{Crash risk.} Carry returns may be steady most of the time and then collapse when carry traders all close their positions at once, as funding dries up. On this view the premium pays for those crashes (Brunnermeier, Nagel and Pedersen 2008).
  \item \textbf{Slow adjustment.} Investors may react slowly to changes in interest rates, so exchange rates move toward UIP only gradually. That would be a mispricing rather than a risk premium, and it should shrink as more money chases it.
\end{itemize}

The three explanations leave different fingerprints. A risk premium should show up as losses in bad times for the world economy. Crash risk should show up as negative skewness, with losses concentrated in a few episodes. Slow adjustment should show up as a premium that shrinks over time and does not depend on market conditions.

\section{The claim}
\label{sec:06-claim}

\begin{callout}
\textbf{Claim (uncovered interest parity).} The currency of a high-interest country is expected to fall by its interest advantage, so borrowing in low-interest currencies and lending in high-interest ones earns nothing on average.
\end{callout}

To test it, three points need to be settled:

\begin{itemize}
  \item \textbf{Over what period.} UIP is a statement about expectations. Over a few months, random exchange-rate movements swamp everything else, so the test needs decades of data and many currencies.
  \item \textbf{Which currencies.} The lab uses nine developed-market currencies against the dollar. Emerging-market currencies offer higher interest gaps but also capital controls, default risk and much higher trading costs.
  \item \textbf{What would count against it.} UIP fails if the Fama slopes are significantly below 1 and a carry portfolio earns a significant average return. It fails more deeply if the return survives trading costs and the slopes are below zero, not merely below one.
\end{itemize}

\section{The original evidence}
\label{sec:06-evidence}

\textbf{Fama (1984).} Using monthly data for nine currencies against the dollar after 1973, Fama found that the slopes were not just below 1 but generally negative: high-interest currencies tended to rise. He showed that this implies a risk premium in forward rates that changes over time and moves more than expected exchange-rate changes do. Froot and Thaler (1990) later summarized the literature: across 75 published estimates, the average slope was about $-0.88$.

\textbf{Carry as a risk factor.} Lustig, Roussanov and Verdelhan (2011) grouped currencies into portfolios by interest rate. The portfolio that was long high-rate and short low-rate currencies\index{factor!HML (value)} earned substantial average returns. A single carry factor, together with a dollar factor, explained much of the difference in returns across currencies, and high-interest currencies were more exposed to the carry factor. Menkhoff, Sarno, Schmeling and Schrimpf (2012) found that carry losses coincided with rises in global currency volatility\index{volatility}.

\textbf{Crashes.} Brunnermeier, Nagel and Pedersen (2008) showed that carry returns are negatively skewed\index{skewness}: large losses are more likely than large gains. The losses came when risk appetite and the availability of funding\index{liquidity} fell, measured among other things by rises in the VIX\index{VIX}, as carry traders were forced to close their positions together.

\textbf{Peso problems.}\index{peso problem|textbf} Another possible explanation is that the high returns pay for rare disasters that have not yet appeared in the data, a \textbf{peso problem}. Burnside, Eichenbaum, Kleshchelski and Rebelo (2011) tested this by buying currency options that protect carry positions against crashes. The protected trade still earned a substantial return. They concluded that rare disasters could explain the premium only if the market prices them far above their likely cost.

\section{Lab~\thechapter: Replicating Fama (1984) and Lustig, Roussanov and Verdelhan (2011)}
\label{sec:06-lab}

\begin{callout}
\textbf{Lab~\ref{ch:06} at a glance.} \emph{Papers replicated:} the Fama (1984) regression; the carry portfolios of Lustig, Roussanov and Verdelhan (2011). \emph{Claim on trial\index{claims-on-trial method}:} high-interest currencies fall enough to cancel their extra yield. \emph{Data:} free monthly data from FRED\index{FRED}: the Federal Reserve's exchange rates for nine currencies (yen, pound, Swiss franc, Canadian, Australian and New Zealand dollars, Swedish krona and Norwegian krone, and the euro from 1999) and OECD three-month interbank rates for each country and the United States. The VIX and the Fama--French\index{Fama--French factors} market factor are optional. \emph{Code:} runs as a script (\texttt{python lab18.py}) or as a notebook with a prediction prompt before each step.
\end{callout}

\textbf{Part A, Fama regressions.} For each currency, the lab regresses the monthly change in the log exchange rate (dollars per unit of foreign currency) on last month's U.S.-minus-foreign interest gap, with Newey--West\index{standard error!Newey--West (HAC)} standard errors. It also estimates a single slope for all currencies together (Lab Figure~\ref{ch:06}.1).

\textbf{Part B, the carry portfolio.} Each month, the lab buys the three currencies with the highest interest rates and sells the three with the lowest, in equal amounts, using rates known at the end of the previous month. It splits the portfolio's return into the interest earned and the change in exchange rates (Lab Figure~\ref{ch:06}.2), and compares the distribution of returns with a normal distribution\index{distribution!normal} (Lab Figure~\ref{ch:06}.3).

\textbf{Part C, robustness checks.} Three subperiods (up to 1998, 1999--2007, and 2008 onward), trading costs\index{transaction costs} of 0 to 20 basis points\index{basis point} per unit traded, exposure to stock returns and to changes in the VIX, and the ten worst months.

\textbf{Part D, the investor's view.} Drawdowns\index{drawdown}, recovery times, the share of losses that came from the worst months, and carry returns plotted against changes in the VIX (Lab Figures~\ref{ch:06}.4 and~\ref{ch:06}.5).

Three choices about the data matter:

\begin{itemize}
  \item The interest gap stands in for the forward discount, which covered interest parity makes nearly equal, because forward rates are not freely available.
  \item FRED's exchange rates are monthly averages of daily rates. Averaging smooths returns and makes crashes look milder than they were.
  \item The sample covers only developed-market currencies, where interest gaps are smaller but trading costs and capital controls are minor.
\end{itemize}

\textbf{What the simulated run shows.} In the simulation, high-interest currencies fall by only 30\% of their interest advantage, and there are rare months of panic in which they fall sharply. The lab finds what it should: a combined Fama slope well below 1; a carry portfolio with a positive average return and negative skewness; interest income partly offset by exchange-rate losses; and losses when the VIX jumps. The slopes for individual currencies are very imprecise, as they are in real data. These results confirm that the code works. They say nothing about real markets. Run the lab on real data to reach your own verdict.

\subsection{Robustness checks}
\label{sec:06-stress}

Results over the full sample can hide weaknesses. The checks below test the claim under the conditions in which it is hardest to meet.

\textbf{Subperiods.} Carry returns have varied a great deal over time. A strategy that earns most of its return in a few long calm stretches and gives much of it back in a few months can look very different depending on where the sample starts and ends. The period since 2008 matters most, because interest rates among developed countries were close to one another for years, leaving little carry to earn.

\textbf{Trading costs.}\index{transaction costs|textbf} The currencies of the largest developed economies are among the most liquid assets in the world, with bid--ask spreads\index{bid--ask spread} of a few basis points. Carry portfolios trade\index{rebalancing} only when the ranking of interest rates changes, so turnover\index{turnover} is modest. Costs reduce the carry premium in developed markets but rarely eliminate it. For emerging-market currencies, they matter much more.

\textbf{Skewness and crashes.}\index{skewness|textbf} Traders describe the carry trade as ``going up the stairs and down the elevator.'' Its returns should be negatively skewed, and its worst months should cluster in episodes of global stress: 1992 (the European exchange-rate crisis), 1998 (the Russian default and the collapse of LTCM), 2008, and the yen rally of August 2024. The lab lists the worst months and their share of total losses.

\textbf{Exposure to stocks and fear.} If carry is a risk premium\index{risk premium}, its returns should rise and fall with stocks, especially in bad months, and fall when the VIX rises. A strategy whose sensitivity to stocks, its beta\index{beta}, rises in falling markets offers less diversification\index{diversification} than its average correlation\index{correlation} suggests.

\textbf{Monthly averages.} Because the exchange rates are monthly averages of daily rates, the lab's worst months understate how fast the real crashes were. In October 2008, several carry currencies fell by more than 10\% against the yen within days.

\subsection*{Your verdict}

Reach your own verdict on the claim from the numbers in your own run, using the scale of Section~\ref{sec:01-method}. Then compare it with the published evidence summarized in Section~\ref{sec:06-verdict}. If your verdict differs, explain which of your results led you there.

\section{Where the evidence stands}
\label{sec:06-verdict}

The published studies cited in this chapter point to a clear pattern: \textbf{high-interest currencies have not fallen by their interest advantage on average, so the carry trade has earned a premium, but that premium is paid for with sudden crashes in global downturns.} The table summarizes where the evidence supports the claim and where it does not, as of November 2026, and names the source of each finding. The lab lets you check these findings with free data.

{\small
\begin{longtable}{@{}>{\raggedright\arraybackslash}p{0.19\linewidth}>{\raggedright\arraybackslash}p{0.10\linewidth}>{\raggedright\arraybackslash}p{\dimexpr 0.48\linewidth-6\tabcolsep\relax}>{\raggedright\arraybackslash}p{0.23\linewidth}@{}}
\toprule
Situation & Holds? & Why & Evidence \\
\midrule
\endfirsthead
\multicolumn{4}{@{}l}{\emph{(continued)}} \\
\toprule
Situation & Holds? & Why & Evidence \\
\midrule
\endhead
\midrule
\multicolumn{4}{r@{}}{\emph{(continued on the next page)}} \\
\endfoot
\bottomrule
\endlastfoot
Fama slope equal to 1 & \textbf{No} & Slopes are below 1 and usually below zero; across 75 published estimates the average was about $-0.88$. & Fama (1984); Froot and Thaler (1990) \\
\addlinespace[3pt]
Carry portfolios earn nothing & \textbf{No} & Portfolios long high-rate and short low-rate currencies earned substantial returns, largely explained by a common carry factor. & Lustig, Roussanov and Verdelhan (2011) \\
\addlinespace[3pt]
Carry in global downturns & \textbf{Yes} & Carry returns are negatively skewed and lose heavily when volatility rises and funding dries up, giving back years of interest income. & Brunnermeier, Nagel and Pedersen (2008); Menkhoff, Sarno, Schmeling and Schrimpf (2012) \\
\addlinespace[3pt]
Rare disasters not yet seen & \textbf{No} & Carry hedged with currency options still earned a substantial return, so unseen disasters explain the premium only if priced far above their likely cost. & Burnside, Eichenbaum, Kleshchelski and Rebelo (2011) \\
\addlinespace[3pt]
Long horizons and very high inflation & \textbf{Partly} & Some studies find slopes closer to 1 at horizons of several years and where large interest gaps are matched by large depreciations. & Later studies; not tested in the lab \\
\addlinespace[3pt]
Carry after 2008 & \textbf{Unclear} & Narrow interest differentials among developed countries shrank the carry available; whether the premium per unit of differential changed is open. & Section~\ref{sec:06-stress} \\
\end{longtable}
}

The premium looks like payment for risk, not like an arbitrage. An investor who earns carry is paid for providing money and insurance to the rest of the world at the moments when both are scarce.

\textbf{What would change the picture.} The evidence would move toward the claim if carry returns since 2008 had been consistently close to zero, as small interest gaps and crowded trades\index{crowding} might imply.

\subsection{Open questions}
\label{sec:06-open}

\begin{itemize}
  \item \textbf{Is carry a risk premium or a mispricing\index{mispricing} that is slowly corrected?} If it were a mispricing, it should shrink as more money chases it. Has it?
  \item \textbf{What happened to carry after 2008?} Interest rates near zero in many developed countries removed most interest gaps. Did the premium per point of interest gap change, or only the opportunity?
  \item \textbf{Do departures from covered interest parity matter?} Since 2008, covered interest parity has not held exactly, reflecting the cost of using banks' balance sheets. How does this affect measured carry returns?
  \item \textbf{Can crashes be avoided?} Rules based on momentum\index{momentum} or volatility claim to sidestep the worst months. Do they work out of sample, or do they simply give up the premium?
\end{itemize}

\section{What practitioners and investors can use}
\label{sec:06-investors}

The evidence in Section~\ref{sec:06-verdict} suggests several practical lessons.

\textbf{What a currency carry fund earns.} A carry strategy collects the differences in interest rates between currencies, and keeps them as long as exchange rates do not move against it. In calm years it delivers modest, steady gains. Its risk is not spread evenly over time. Most of its losses come in a few months of global stress, when high-yielding currencies fall together and the borrowed currencies, above all the yen and the Swiss franc, rise.

\textbf{How it loses money.} Carry crashes are fast. They begin when something forces carry traders who have borrowed heavily to cut their positions: a rise in volatility, a margin call\index{margin}, or a surprise from the central bank of a borrowed currency. Their selling pushes high-yielding\index{credit rating} currencies down further, which forces more selling. Spreading the trade across many currency pairs helps little, because the losses come from a common source.

\textbf{Carry in a portfolio.} Because carry tends to lose money when stocks fall, it diversifies a stock portfolio less than its average correlation with stocks suggests. Investors should measure that correlation in bad months, not over the whole sample.

\textbf{Questions to ask a manager running currency carry:}

\begin{itemize}
  \item Which currencies does the fund borrow in, and how much of the portfolio depends on one or two currency pairs?
  \item How much leverage\index{leverage} does the fund use, and what happens to its margin in a volatility spike?
  \item How did the strategy perform in late 2008 and in August 2024?
  \item Does it use a rule to reduce exposure when volatility rises, and how did that rule perform out of sample\index{in-sample versus out-of-sample}?
\end{itemize}

\begin{callout}
\textbf{What Failed: the yen carry trade, 2008.} In the years before 2008, borrowing yen at interest rates near zero and investing in Australian dollars, New Zealand dollars and other high-yielding currencies was one of the most popular trades in the world, among Japanese retail investors and hedge funds alike. When the financial crisis deepened in the autumn of 2008, carry traders rushed to repay their yen loans at the same time. Between July and late October 2008, the Australian and New Zealand dollars lost more than a third of their value against the yen, wiping out years of interest income in a few weeks.

The episode matched the risk-premium explanation exactly. The trade paid steadily in good times and lost heavily in the worst global downturn in decades, when investors could least afford it. It happened again on a smaller scale in August 2024. After the Bank of Japan raised its policy rate to 0.25\% at the end of July, the yen rose sharply, carry positions were cut, and on August 5 the Nikkei 225 fell 12.4\%.
\end{callout}

\section*{Exercises}
\addcontentsline{toc}{section}{Exercises}

\subsection*{Conceptual (all readers)}

\begin{enumerate}
  \item The U.S. one-year rate is 4\% and the Japanese rate is 0.5\%. The spot rate is 150 yen per dollar. Compute the one-year forward rate. Under UIP, what spot rate is expected in a year?
  \item Using the numbers in Exercise 1, compute the return to a dollar-based investor who borrows yen, converts them to dollars and invests at 4\%, if the spot rate stays at 150, moves to the UIP forecast, or falls to 135.
  \item Explain why covered interest parity is enforced by arbitrage but uncovered interest parity is not.
\end{enumerate}

\subsection*{Lab extensions (technical readers)}

\begin{enumerate}
  \setcounter{enumi}{3}
  \item Part A: how many currencies have Fama slopes significantly below 1? How many below 0? What is the slope for all currencies together?
  \item Part B: what share of the carry portfolio's interest income was lost to exchange-rate changes? Compare it with what UIP predicts.
  \item Part C: compare carry returns before and after 2008. What happened to the interest gaps?
  \item Part C: estimate the carry portfolio's beta to stocks in all months and in months when stocks fell. What does the difference mean for an investor?
  \item Add a rule that holds the carry portfolio only when the VIX is below its median over the previous 12 months. Compare the Sharpe ratio and the worst month with those of the strategy without the rule.
  \item A financial adviser recommends a fund that earns ``the interest-rate differential between currencies.'' Using this chapter's evidence, explain what the fund is really paid for and when it is likely to lose money.
\end{enumerate}
